\chapter{Cognitive Architecture}

The intelligence layer requires a cognitive architecture: an explicit organisation of how
the model makes decisions over time. In early agentic systems this organisation was often
implicit, embedded in a single long prompt. In production systems it must be explicit,
because it is the cognitive architecture that determines which model to invoke, how much
reasoning to allocate to a given task, whether to retrieve additional context or act on
what is already known, whether to seek clarification from the user, and when to stop.

\section{The Agent Loop}

The defining property of an agentic system is that it pursues a goal across multiple
steps rather than producing a single response. This iterative structure is what the
cognitive architecture governs, and it is what distinguishes an agent from a prompted
language model.

The loop is necessary because the tasks assigned to agents are rarely solvable in a
single pass. A software engineering task may require reading failing tests, forming a
reasoning trace about the probable cause, modifying code, running tests again, observing
a new failure, revising the reasoning, and repeating. A research task may require
querying multiple sources, discovering that initial results are insufficient, broadening
the search strategy, and synthesising across an expanded set of evidence. In each case
the agent must maintain a coherent picture of where it is in the task, what it has
learned, and what it should attempt next.

The standard formulation of this loop, introduced by the ReAct framework (Yao et al.,
2022), interleaves reasoning and acting in three repeating stages. In the thought stage,
the model generates an internal reasoning trace: it analyses the current state of the
task and determines what action is most likely to advance it. In the action stage, it
executes that action---calling a tool, retrieving context, producing a partial answer,
or requesting clarification. In the observation stage, it receives the result of the
action and incorporates it into its reasoning before the next thought. This
thought-action-observation cycle continues until a termination condition is met. The
critical property of the loop is that the agent learns from each observation before
deciding on the next step: a tool call that returns an error is an observation that
informs the subsequent thought, rather than a terminal failure.

\section{State and the Reasoning Trace}

For the loop to function across multiple iterations, the agent must maintain state.
State is the accumulated record of what the agent has done and learned: the user's
original input, the message history, retrieved documents, tool outputs, the current
plan, completed steps, unresolved questions, and accumulated cost.

The most important architectural distinction within state is between observations and
the reasoning trace. An observation is a fact returned by an external source---a tool
output, a retrieved document, a user message, a test result. The reasoning trace is the
agent's own internal thought: a hypothesis about why a test is failing, a candidate plan
for the next steps, an interpretation of an ambiguous user request. Conflating these two
categories is a primary source of hallucination and reasoning failure. An agent that
treats its own reasoning trace as though it were an observed fact will build subsequent
steps on an unverified foundation and propagate the error forward through the loop.

A sound architectural practice is to maintain a clear separation between what the agent
has observed and what it has inferred. The final response should be grounded in
observations---facts returned by external sources---rather than in intermediate reasoning
steps that were never confirmed by a subsequent observation. This discipline is what
allows an agent to recover from an incorrect initial hypothesis: when a reasoning trace
leads to an action whose observation contradicts the hypothesis, the agent can revise its
reasoning and converge on a correct solution rather than compounding the error.

\section{Deliberation Budgets}

Reasoning is a resource, and a well-designed cognitive architecture makes its allocation
explicit. A deliberation budget specifies, for a given task or task class, the model to
be used, the maximum number of output tokens, the permitted number of tool iterations,
the number of retrieval rounds, and whether the system employs self-checking or
verification passes. These parameters collectively determine both the cost and the
quality ceiling of the agent's response.

The budget should be visible in the execution trace. A task that consumes substantially
more deliberation than its complexity warrants represents a cost failure; a task whose
complexity exceeds the deliberation allocated to it represents a quality failure. Making
the budget explicit transforms these from invisible degradations into diagnosable
conditions.

\section{Clarification and Autonomy}

A recurring decision in the agent loop is whether to seek clarification from the user
or to proceed autonomously. The appropriate choice depends on the risk profile of the
next action and the recoverability of an error.

Clarification is warranted when required parameters are absent, when multiple plausible
interpretations of the user's intent would lead to materially different actions, when
the next action is irreversible, when the cost of an incorrect action is high, or when
the applicable permissions are ambiguous. Autonomous continuation is warranted when the
next step carries low risk, when the agent can verify the result of its action before
committing to it, when the operation is reversible, when the user has explicitly
delegated the task, or when missing information can be retrieved safely without user
involvement. The design of this decision boundary is a governance choice as much as a
technical one: it determines the degree of autonomy the agent exercises and the degree
of accountability it shares with the user.

\section{Stopping Conditions}

Because the loop has no natural end when the task is open-ended, the cognitive
architecture must specify explicit termination conditions. Without them the agent
continues iterating past the point at which the task is complete, consuming budget and
potentially degrading a solution that was already correct.

Termination conditions fall into four classes. Success-based termination occurs when the
required output has been produced and verified---the test passes, the requested document
has been retrieved and confirmed relevant, the generated artefact meets the specified
criteria. Budget-based termination occurs when the agent has exhausted its allocation of
tool calls, tokens, time, or cost. Safety-based termination occurs when the next action
would require permissions that the user or policy has withheld. Uncertainty-based
termination occurs when the agent cannot make further progress autonomously and a human
decision is required to proceed.

A property of this architecture that deserves explicit attention is that success-based
termination is itself a model judgment, rather than a deterministic rule. The agent uses
the language model to assess whether the current state of the task satisfies the original
goal. This means the stopping decision is subject to the same reasoning capabilities and
failure modes as any other step in the loop: the model may judge a task complete when it
is not, or continue iterating when sufficient evidence already exists. Budget-based
termination is the only class that is enforced deterministically by the framework rather
than decided by the model. This asymmetry---model-judged success, framework-enforced
budget---is a structural property of the cognitive architecture that the evaluation layer
must account for, by measuring not only whether the agent reached a correct answer but
whether it stopped at the right point.

\section{Summary}

The cognitive architecture controls how the model reasons and acts over time. It
structures the agent's execution as a thought-action-observation loop, allocates
deliberation budgets, maintains a clear separation between observations and the
reasoning trace, governs the boundary between clarification and autonomous action, and
enforces termination conditions. The prompt and skill content that informs each reasoning
step belongs to the memory and knowledge layer; the cognitive architecture determines
how and when that content is used, and for how long the agent continues to pursue the
goal.